\section{Despliegue con Docker Compose}
\label{sec:docker}

El despliegue con Docker elimina la necesidad de instalar manualmente GStreamer, Python
y sus dependencias en el equipo anfitrión. Todo el sistema se ejecuta dentro de
contenedores reproducibles, de modo que el comportamiento es idéntico en cualquier
máquina con Docker instalado. El despliegue se define en dos ficheros:
\texttt{Dockerfile} (imagen base del sistema) y \texttt{docker-compose.yml}
(conjunto de servicios).

El \texttt{Dockerfile} original del repositorio (año 2016) usaba Ubuntu~15.10 con
dependencias ya obsoletas e imágenes de Docker Hub que ya no existen. Se reescribió
desde cero sobre \texttt{ubuntu:22.04 LTS}.

\subsection{Dockerfile}

La imagen se construye sobre Ubuntu~22.04~LTS. Todas las dependencias se instalan en
una sola capa para reducir el tamaño final:

\begin{itemize}
  \item GStreamer~1.0 con sus plugins \textit{base}, \textit{good}, \textit{bad} y
        \textit{ugly}.
  \item FFmpeg, para la inyección de fuentes de vídeo.
  \item \textit{Bindings} Python de GObject (\texttt{python3-gi}), la biblioteca Cairo
        para renderizado de \textit{overlays}, y la librería AMQP (\texttt{pika}) para
        la integración con RabbitMQ.
  \item Utilidades de red (\texttt{netcat}, \texttt{iproute2}) para el
        \textit{healthcheck} y el autodescubrimiento de IP.
  \item Bibliotecas matemáticas \texttt{scipy} y \texttt{numpy} para el cálculo de
        curvas B-Spline en las transiciones animadas.
\end{itemize}

Un \textit{healthcheck} basado en \texttt{nc -zw1 localhost 9999} permite a Docker
Compose saber cuándo voctocore está listo antes de lanzar los servicios que dependen
de él.

\subsection{Docker Compose y espacio de red compartido}
\label{sec:docker_compose}

Voctomix requiere que las fuentes de cámara se conecten al núcleo mediante
\texttt{localhost}. En Docker, cada contenedor tiene su propio \texttt{localhost}
aislado, lo que rompe todas las conexiones TCP internas.

La solución adoptada es el patrón \textbf{\textit{Shared Network Namespace}}: los
contenedores de cámara y el servicio de telemetría declaran
\texttt{network\_mode: "service:voctocore"}, lo que hace que compartan la misma
interfaz de red que el contenedor del núcleo. Desde el punto de vista del sistema
operativo, todos estos procesos ven el mismo \texttt{localhost}, sin necesidad de
modificar ninguna línea de código.

El fichero \texttt{docker-compose.yml} define los siguientes servicios:

\begin{table}[H]
  \centering
  \caption{Servicios definidos en \texttt{docker-compose.yml}.}
  \label{tab:servicios_docker}
  \small
  \begin{tabular}{|l|l|c|l|}
    \hline
    \textbf{Servicio} & \textbf{Imagen} & \textbf{Puertos expuestos} & \textbf{Función} \\
    \hline
    voctocore       & local/voctomix & 9999, 11000--12000, 14000--14005, 15000 & Núcleo multimedia \\
    cam1--cam4      & local/voctomix & (namespace compartido) & Fuentes de prueba FFmpeg \\
    stream\_blanker & local/voctomix & 17000--17001, 18000 & Bucles pausa/\textit{offline} \\
    rabbitmq        & rabbitmq:mgmt  & 5672, 15672 & Broker AMQP \\
    telemetria      & local/voctomix & 8080 & Telemetría REST/AMQP \\
    \hline
  \end{tabular}
\end{table}

\subsection{GUI nativa y volúmenes}

La interfaz gráfica voctogui se ejecuta de forma nativa en el sistema anfitrión y se
conecta al núcleo Docker a través del puerto 9999 expuesto. Virtualizar una interfaz
GTK dentro de un contenedor introduce problemas de latencia y aceleración por
hardware (GPU/VAAPI) que hacen inviable esa opción en producción.

Los logos e imágenes de fondo que GStreamer necesita se inyectan en el contenedor
mediante un volumen \textit{bind mount}: \texttt{./images:/opt/images}. Esto evita
modificar el código fuente para adaptar las rutas a la estructura interna del
contenedor.

\subsection{Resiliencia al arranque}

Durante las pruebas se observó que los contenedores de cámara arrancaban antes de
que voctocore abriera los puertos TCP, produciendo fallos de conexión silenciosos. El
problema se resolvió con dos mecanismos:

\begin{itemize}
  \item \textbf{\texttt{restart: always}:} Docker reinicia automáticamente cualquier
        contenedor que falle, sin intervención del operador.
  \item \textbf{Esperas activas:} en los scripts de arranque nativos, los
        \texttt{sleep} fijos se sustituyeron por bucles que comprueban con
        \texttt{ss -lnt} que el puerto 9999 está en escucha antes de lanzar los
        servicios dependientes.
\end{itemize}

Además, la opción \texttt{SO\_REUSEADDR} en el servidor HTTP de telemetría evita el
error \texttt{Errno 98} que bloqueaba el puerto 8080 tras reinicios rápidos.
